\phantomsection
\chapter*{Introducció}
\addcontentsline{toc}{chapter}{Introducció}
\label{cap:introduccio}

Aquest llibre desenvolupa una teoria estructural de conjunts en la qual les aplicacions són els únics individus primitius. Els conjunts, els elements i les construccions habituals es recuperen progressivament a partir dels dominis, els codominis i la composició.


\phantomsection
\section*{Una mirada diferent: el punt de vista estructural}
\addcontentsline{toc}{section}{Una mirada diferent: el punt de vista estructural}

En la teoria clàssica de conjunts ---la formulació tradicional de
Zermelo--Fraenkel--- els protagonistes indiscutibles són els
\textbf{elements}. Un conjunt es descriu visualment com un «sac» o una
«col·lecció» on guardem objectes dins, i la relació fonamental és la pertinença \( \in \), mitjançant la qual diem que un element \(x\) pertany al conjunt \(A\) escrivint \(x \in A\). Les aplicacions són actors secundaris, definides a posteriori com a conjunts de parells ordenats.

En aquest llibre canviarem la perspectiva: \textbf{els
protagonistes seran les aplicacions}.

Farem com en geometria: per entendre les figures sovint interessa més
saber \textbf{com es transformen o es relacionen entre elles} ---mitjançant
moviments, projeccions o simetries--- que no pas preguntar-se de què estan
fetes per dins. Adoptarem la mateixa actitud amb els conjunts: ens
interessarà com es relacionen mitjançant les aplicacions, més que no pas
la seva estructura interna de pertinença.

\begin{center}
\begin{tikzpicture}[scale=1,every node/.style={font=\small}]

%--------------------------------------------------
% Visió clàssica
%--------------------------------------------------

\node[font=\bfseries] at (1.5,3.2) {Visió clàssica (ZF)};

\draw[thick,rounded corners=14pt] (0,0) rectangle (3,2.5);

\node at (1.5,2.75) {$A$};

\fill (0.8,1.8) circle (2pt) node[right=2pt] {$x$};
\fill (2.2,1.0) circle (2pt) node[right=2pt] {$y$};
\fill (1.1,0.5) circle (2pt) node[right=2pt] {$z$};

\node at (1.5,-0.45)
{Contenidor d'elements: $x\in A$};

%--------------------------------------------------
% Fletxa central
%--------------------------------------------------

\draw[-{Stealth[scale=1.4]},line width=1.2pt]
(4.0,1.25)--(5.8,1.25);

%--------------------------------------------------
% Visió estructural
%--------------------------------------------------

\node[font=\bfseries] at (8.5,3.2)
{Visió estructural (aquest llibre)};

\node (A) at (7.2,1.8) {$A$};
\node (B) at (9.6,1.8) {$B$};
\node (C) at (8.4,0.4) {$C$};

% Aplicacions

\draw[-{Stealth[scale=1.2]},thick]
(A)--node[above]{$f$}(B);

\draw[-{Stealth[scale=1.2]},thick]
(A)--node[left]{$g$}(C);

\draw[-{Stealth[scale=1.2]},thick]
(C)--node[right]{$h$}(B);

%--------------------------------------------------
% Identitats
%--------------------------------------------------

% id_A
\draw[-{Stealth[scale=1]},thick]
(7.20,2.00)
.. controls (6.80,2.45) and (7.60,2.45) ..
(7.20,2.00);
\node[font=\scriptsize] at (6.55,2.48)
{$\id_A$};

% id_B
\draw[-{Stealth[scale=1]},thick]
(9.60,2.00)
.. controls (9.20,2.45) and (10.00,2.45) ..
(9.60,2.00);
\node[font=\scriptsize] at (10.25,2.48)
{$\id_B$};

% id_C
\draw[-{Stealth[scale=1]},thick]
(8.40,0.20)
.. controls (8.00,-0.25) and (8.80,-0.25) ..
(8.40,0.20);
\node[font=\scriptsize] at (9.05,-0.28)
{$\id_C$};

\node at (8.4,-0.85)
{Xarxa d'aplicacions};

\end{tikzpicture}
\end{center}

\subsection*{Una mica de context: d'on ve aquesta idea?}

Aquesta manera de reconsiderar els fonaments de la matemàtica \textbf{no és
nova}. La va proposar el matemàtic F.~William Lawvere durant la dècada de
1960~\cite{lawvere1964}, en adonar-se que la pràctica quotidiana ---en àlgebra, topologia o
anàlisi--- consisteix a estudiar com els objectes interactuen entre ells
mitjançant aplicacions, i no pas a registrar la naturalesa dels seus
elements interns.

Lawvere ho va formular en el llenguatge de la teoria de categories, en el
que es coneix com la \emph{teoria estructural de conjunts de Lawvere}
(ETCS)~\cite{lawvere1964,lawvere-etcs-reprint}, però \textbf{per seguir aquest llibre no cal conèixer la teoria de categories}. Una presentació moderna, propera a l'esperit d'aquest llibre, es troba a \cite{lawvere-rosebrugh}; per a la lectura des de la teoria de topos, vegeu \cite{maclane-moerdijk}.

El nostre objectiu és construir una teoria estructural de conjunts des de zero, amb un llenguatge formal de primer ordre elemental i amb els convenis habituals que permeten abreujar la notació.


\section*{Criteri d'escriptura formal}\label{not:criteri-formal}
Les definicions i els axiomes s'escriuran mitjançant fórmules del llenguatge formal o mitjançant abreviacions introduïdes explícitament. En les proposicions, els teoremes i les demostracions farem servir sistemàticament aquestes abreviacions, perquè la notació primitiva no amagui l'estructura matemàtica dels arguments.

Cada conveni introduït en un entorn de notació és eliminable: qualsevol expressió que el contingui es pot traduir a una fórmula construïda amb \(\Dom\), \(\Cod\), \(\Com\), la igualtat i els predicats definits anteriorment. Per tant, aquests convenis no amplien la teoria; només en simplifiquen l'escriptura.

La presentació alterna tres nivells. El nivell formal fixa les definicions i els axiomes; el nivell geomètric els representa mitjançant diagrames externs de fletxes; i les observacions de lectura clàssica expliquen com es recuperen els conjunts i les funcions habituals. Aquests tres nivells es mantindran diferenciats al llarg de l'obra.
